% 附录 C 原文勘误与存疑（汇总各译手的逐页对勘记录，定稿前须与正文逐条复核）
\section{原文勘误与存疑}

本附录登记译文与原文逐页对勘时发现的笔误、排印错误与存疑之处。凡属笔误，正文一般照原文译出，
另在译注中说明；本表只作索引，便于读者复核。页码一律为原文页码（对照见附录 B）。
表中「处理」一栏说明译文如何落笔。

\subsection*{一、原文笔误（已按数学意义修正，并在正文加注）}

\begin{longtable}{@{}p{0.10\textwidth}p{0.18\textwidth}p{0.65\textwidth}@{}}
\toprule
\textbf{原文页} & \textbf{位置} & \textbf{说明与处理} \\
\midrule
\endhead
\bottomrule
\endfoot
1 & \S1 第 4 段 & 原文「If $[\omega^{\vee}]$ does not belong to $H^2(X,\Z)$」中 $X$ 应为 $X^{\vee}$：
讨论对象是辛侧流形 $X^{\vee}$ 的辛类。译文按 $X^{\vee}$ 译出，并加译注。\\
9 & \S2.5 & 足标自相矛盾：同一处 $SU(2)$ 的 WZW 模型中心荷以 $\phi_n$ 与 $\phi_k$ 两种记号混写。
译文照原文译出，另加注说明应为一种记号。\\
11 & \S2.7 & 讲镜像对称等同某一邻域的那句话多出一个右括号。译文按原意（一次性括起两个邻域）处理，不影响数学。\\
14 & \S3.1，命题 1 & 同一命题内 $T-\mathrm{Id}$ 与 $T-\mathrm{id}$ 混用。译文统一记为 $T-\Id$。\\
25 & \S4.3，定义 7 & 原文作 $m_j(a_l,\dots,a_{l+j}),a_{l+j+1},\dots$，与 $m_j$ 取 $j$ 个变元、以及
p23 定义 3 的同一方程都不合。译文改为 $m_j(a_l,\dots,a_{l+j-1}),a_{l+j},\dots$。\\
26 & \S4.3 首段 & 原文「differs from the case of non-unital $A_\infty$-pre-categories」；此处比较对象是
定义 6（范畴的函子）与定义 8（预范畴的函子）所诱导的 $A_\infty$-态射，故应为「非幺 $A_\infty$-范畴」。\\
27 & \S4.3，定义 13 & 原文作「$C_0=C$，$C_0=D$」与「$0\le i\le k-1$」，与所取有限序列 $(C_0,\dots,C_n)$ 不合。
译文改为 $C_n=D$ 与 $0\le i\le n-1$。\\
46 & \S6.5，引理 2 & 三处排印疏漏：第 2) 款 $Z_1,Z_2\in Y$ 应为 $\subset Y$；等价链末项多一个右括号；
末段第二个 $m_k^{M(Y)}$ 应为 $m_k^{DR_0^{tr,\Pi}(Y)}$。\\
46 & \S6.6 末，定理 3 之后 & 「move on $Y$ with with different speeds」重出一个 with。纯英文笔误，中文不体现。\\
57, 58 & \S8.3 & 把 $C^{tr,\Pi}_{unram,0}(Y)$ 拼作 $C^{tr,\Pi}_{unrum,0}(Y)$（两处）。译文统一为 unram。\\
\end{longtable}

\subsection*{二、存疑与待考（照原文译出，未擅改）}

\begin{longtable}{@{}p{0.10\textwidth}p{0.18\textwidth}p{0.65\textwidth}@{}}
\toprule
\textbf{原文页} & \textbf{位置} & \textbf{说明} \\
\midrule
\endhead
\bottomrule
\endfoot
5 & \S2.1 & 中心荷的渐近式与常见的 Cardy 公式差一个 $\sqrt{2}$ 因子。原文如此，译文照录，未加注以免越位。\\
13 & \S3.1，定义 1 之前 & 原文对 $X^{\mathrm{mer}}_q$ 与 $X_q$ 的使用时有出入，读者宜按上下文区分
「$q=0$ 处的形式邻域上的模型」与「族中的一般纤维」。\\
17 & \S3.2，定义 2 & Monge--Ampère 流形的三元组 $(Y,g,\nabla)$ 中，原文未明说 $\nabla$ 的平坦性；
由后文命题 2 与推论 1 可知 $\nabla$ 是平坦的仿射结构联络。\\
21 & \S3.3 & 「$Y$ 的分片线性闭子集」：原文作 $Y$ 而非 $\bar Y$；若按「闭包在 $\bar Y$ 中」理解亦通。\\
40 & \S6.3，定理 2 之前 & 原文排印有一处公式错位（文字提取层呈现为孤立的「$1\ X$」），
已据页面图像还原为 $w_{\gamma}=\exp\bigl(-\tfrac{1}{\eps}\sum_{e\in E(T)}\mathrm{var}_{\gamma}(f_{l(e)}-f_{r(e)})\bigr)$。\\
50 & \S7.2，命题 9 & 陈述中两侧均写作 $Y^{\mathrm{an}}$，疑原文有笔误（或其中一侧应带某种完备化）。
译文照原文，另加译注说明。\\
58 & \S8.3，命题 14 & 第二个公式用 $\Hom(X_i,X_{i+1})$（对象作自变量），第一个公式用 $\Hom(E_i,E_{i+1})$（层作自变量），
二者不统一；系原文写法，译文照录。\\
\end{longtable}

\subsection*{三、译文的排印改动（不涉及数学内容）}

\begin{itemize}[leftmargin=1.8em,itemsep=2pt]
\item 原文一律不编号的行间公式，译文亦不编号；原文的行内长公式在译文中偶改为行间公式，符号一个未改。
\item 原文三幅手绘示意图（p10、p11、p12）无图题、无编号；译文直接裁切原文图版（600 dpi 重采样），
不加图题编号，也不重绘——线描细节（A/T/B 标号、角点圆弧、箭头）保持原样。
\item 原文脚注 6 条全部保留；译文自己的说明一律入橙色「译注」框，不与原文脚注混排。
\item 原文参考文献（p63--66，34 条）照原样保留英文，未译。
\end{itemize}
